\chapter{The Crypto Trading Business}\label{ch:m3:the-crypto-trading-business}

A \omterm{def:m3:blockchains-for-traders:key}{token} project about to list on exchanges needs someone to quote it. It lends a market maker two percent
of its \omterm{def:m3:blockchains-for-traders:key}{token} supply for a year and grants it call options on the same amount, and the market maker agrees
to keep a bid and an offer on several venues. Neither side pays cash; both think they have a good deal. In
October 2024 the US securities regulator charged three firms calling themselves market
makers with selling, to \omterm{def:m3:blockchains-for-traders:key}{token} promoters, trading that was not trading at all: volume printed between their
own accounts to make \omterm{def:m3:blockchains-for-traders:key}{tokens} look liquid. This chapter is about the business around the markets of the
previous chapters: the agreements under which firms make markets in new \omterm{def:m3:blockchains-for-traders:key}{tokens} and how to value them,
the over-the-counter desks that trade in size, the custody and treasury operations that keep the coins
safe, the operations of a market that never closes, and the regulation that is now arriving.

\section{Token market-making agreements}

\begin{definition}[Token market-making agreement, loan-plus-call structure]\label{def:m3:the-crypto-trading-business:tmma}
A \emph{token market-making agreement}\index{token market-making agreement} is a contract in which a \omterm{def:m3:blockchains-for-traders:key}{token}'s issuer engages
a firm to provide quotes in the \omterm{def:m3:blockchains-for-traders:key}{token} on stated venues, with obligations on spread, depth and presence. In a
\emph{loan-plus-call structure}\index{loan-plus-call structure} the issuer pays for the service by lending the firm \omterm{def:m3:blockchains-for-traders:key}{tokens}
for the term, to be returned at the end, and granting it European call options on \omterm{def:m3:blockchains-for-traders:key}{tokens} at agreed strikes.
\end{definition}

The loan gives the market maker the inventory it needs to offer \omterm{def:m3:blockchains-for-traders:key}{tokens} it does not own; the calls are its pay. The
loan by itself carries no net exposure for the market maker: it holds the \omterm{def:m3:blockchains-for-traders:key}{tokens} and owes the same number back. The calls
make it long the \omterm{def:m3:blockchains-for-traders:key}{token}. A market maker that wants to earn the calls' value rather than bet on the \omterm{def:m3:blockchains-for-traders:key}{token} sells, on day one,
as many \omterm{def:m3:blockchains-for-traders:key}{tokens} as the calls' delta (One Quant Book 1, chapter 25), and keeps that hedge adjusted as the price moves: it
sells as the price rises and buys as it falls. The hedge is sold into the market that has just opened, from the lent
\omterm{def:m3:blockchains-for-traders:key}{tokens}.

\begin{proposition}[Value and hedge of a loan-plus-call deal]\label{prop:m3:the-crypto-trading-business:deal}
Let the issuer lend $N$ \omterm{def:m3:blockchains-for-traders:key}{tokens} at price $S$ for $T$ years and grant calls on $n_i$ \omterm{def:m3:blockchains-for-traders:key}{tokens} at strikes $K_i$. With volatility
$\sigma$ and no drift or discounting, the package is worth $\sum_i n_i C(S, K_i, T, \sigma)$, where $C$ is Black's call price,
and its value divided by $NS$ is the fee the deal implies. The market maker is delta-neutral after selling
$\sum_i n_i \Phi(d_{1,i})$ \omterm{def:m3:blockchains-for-traders:key}{tokens}, with $d_{1,i} = \ln(S/K_i)/(\sigma\sqrt T) + \tfrac12\sigma\sqrt T$.
\end{proposition}

\begin{proof}
\omterm{def:m3:blockchains-for-traders:key}{Tokens} received and \omterm{def:m3:blockchains-for-traders:key}{tokens} owed cancel, so the only exposure is the calls'. Their value and delta under a driftless
lognormal price are Black's formula and $\Phi(d_1)$ per \omterm{def:m3:blockchains-for-traders:key}{token} (One Quant Book 4 for the Monte Carlo version the tutorial uses).
\end{proof}

\begin{omfigure}
\centering
\begin{tikzpicture}[font=\small,
  p/.style={draw,thick,rounded corners=2pt,align=center,minimum height=1cm,text width=2.8cm},
  arr/.style={-{Stealth[length=2.2mm]},thick}]
\node[p,draw=omDef,fill=omDef!8,text width=2.2cm] (iss) at (0,0) {token issuer};
\node[p,draw=omThm,fill=omThm!8,text width=2.2cm] (mm) at (5.4,0) {market maker};
\node[p,draw=omProp,fill=omProp!8,text width=2.2cm] (vx) at (10.2,0) {exchanges};
\draw[arr] ([yshift=8pt]iss.east) -- node[above,font=\footnotesize,align=center] {loan of $N$ tokens,\\call options} ([yshift=8pt]mm.west);
\draw[arr] ([yshift=-8pt]mm.west) -- node[below,font=\footnotesize,align=center] {quotes; $N$ tokens\\back at the end} ([yshift=-8pt]iss.east);
\draw[arr,{Stealth[length=2.2mm]}-{Stealth[length=2.2mm]}] (mm) -- node[above,font=\footnotesize] {bids, offers} node[below,font=\footnotesize] {delta hedge} (vx);
\end{tikzpicture}

\omcaption{A loan-plus-call deal. The issuer pays with a loan of \omterm{def:m3:blockchains-for-traders:key}{tokens} and call options; the market maker quotes on the exchanges and hedges the
calls' delta by selling part of the lent \omterm{def:m3:blockchains-for-traders:key}{tokens}. Schematic.}\label{fig:m3:the-crypto-trading-business:deal}
\end{omfigure}

The structure has a conflict built in. A market maker long calls gains if the price rises above the strikes, and one that is
paid in \omterm{def:m3:blockchains-for-traders:key}{tokens} rather than cash may be tempted to push it there; an issuer that wants a high listing price and a volume ranking may
want the same. Obligations measured on spread and depth can be met honestly; obligations measured on volume invite the \omterm{def:m3:spot-markets:wash}{wash trading}
of \cref{ch:m3:spot-markets}.

\begin{dated}{2026-09}{Market-manipulation-as-a-service}\label{dat:m3:the-crypto-trading-business:mirrors}
On 9 October 2024 the SEC announced fraud charges against three companies purporting to be market makers (ZM Quant, Gotbit and CLS Global) and nine
individuals, alleging that \omterm{def:m3:blockchains-for-traders:key}{token} promoters hired ZM Quant and Gotbit to provide ``market-manipulation-as-a-service'', including artificial trading
volume, for \omterm{def:m3:blockchains-for-traders:key}{tokens} offered to retail investors, and that ZM Quant and CLS Global did the same for a \omterm{def:m3:blockchains-for-traders:key}{token} created at the direction of the FBI as part
of a parallel investigation. The same day the US Attorney's Office in Boston unsealed criminal charges against eighteen individuals and entities,
including four market makers, seized more than USD~25 million in crypto and deactivated trading bots that had wash traded about 60 \omterm{def:m3:blockchains-for-traders:key}{tokens}.
\end{dated}

\begin{definition}[Token unlock]\label{def:m3:the-crypto-trading-business:unlock}
A \emph{token unlock}\index{token unlock} is the scheduled release, from a lock-up, of \omterm{def:m3:blockchains-for-traders:key}{tokens} allocated to a project's team, investors or treasury, after
which their holders may sell them.
\end{definition}

Unlocks add supply at known dates, and the market maker's hedge and its quotes must allow for them: a large unlock is a known future
seller, as an index deletion is in One Quant Book 1, chapter 15.

\begin{omfigure}
\begin{tikzpicture}
\begin{axis}[width=11cm,height=5cm,
  xlabel={token volatility, \% a year},ylabel={implied fee, \% of loan},
  tick label style={font=\small},label style={font=\small},
  xmin=40,xmax=240,ymin=0,ymax=70,grid=major,grid style={gray!25},
  table/col sep=comma]
\addplot[omThm,very thick] table[x=vol,y=fee]{figdata/markets-3/24-the-crypto-trading-business/fee.csv};
\end{axis}
\end{tikzpicture}

\omcaption{The fee implied by the tutorial's listing deal (2\% of supply lent for a year at a listing price of USD~0.50, calls on a third each at 0.75,
1.00 and 1.50) as a function of the \omterm{def:m3:blockchains-for-traders:key}{token}'s volatility (\cref{prop:m3:the-crypto-trading-business:deal}). At 120\% the calls are worth 26\% of the loan.
Illustrative deal. Data: the chapter's tutorial.}\label{fig:m3:the-crypto-trading-business:fee}
\end{omfigure}

\begin{omfigure}
\begin{tikzpicture}
\begin{axis}[width=11cm,height=5cm,
  xlabel={token price, USD},ylabel={hedge, million tokens short},
  tick label style={font=\small},label style={font=\small},
  xmin=0.1,xmax=3,ymin=0,ymax=2.1,grid=major,grid style={gray!25},
  legend columns=2,legend style={font=\footnotesize,at={(0.5,-0.3)},anchor=north,draw=none,fill=none,/tikz/every even column/.append style={column sep=8pt}},
  table/col sep=comma]
\addplot[omThm,very thick] table[x=price,y=t1]{figdata/markets-3/24-the-crypto-trading-business/hedge.csv};\addlegendentry{at signing (a year left)}
\addplot[omDef,very thick,dashed] table[x=price,y=t05]{figdata/markets-3/24-the-crypto-trading-business/hedge.csv};\addlegendentry{half a year left}
\end{axis}
\end{tikzpicture}

\omcaption{The market maker's delta hedge on the listing deal: \omterm{def:m3:blockchains-for-traders:key}{tokens} to hold short against the calls as a function of the \omterm{def:m3:blockchains-for-traders:key}{token} price. It sells as the
price rises (towards the 2 million \omterm{def:m3:blockchains-for-traders:key}{tokens} of the calls) and buys back as it falls. Data: the chapter's tutorial.}\label{fig:m3:the-crypto-trading-business:hedge}
\end{omfigure}

\section{OTC desks}

Large trades in crypto are done off exchange, by request for quote (One Quant Book 2, chapter 22): a fund asks a desk for a price in 500 bitcoin, the desk
quotes and, if hit, hedges on exchanges over minutes or hours. The desk's price is the hedge's expected cost (the depth it will consume across venues, fees,
the \omterm{def:m3:blockchains-for-traders:stable}{stablecoin} and fiat legs) plus the risk of the price moving while it hedges, plus a margin. The settlement is the difficult part: fiat moves through banks
in banking hours and coins on chains at any hour, so one side of the trade is always in flight. Desks settle against pre-funded balances, or deliver first
to clients with credit lines, and the credit risk of the settlement window is priced like the settlement risk of the currency markets (One Quant Book 2, chapter 20).

\section{Treasury and custody operations}

\begin{definition}[Self-custody, hot wallet, cold wallet, multi-party computation]\label{def:m3:the-crypto-trading-business:custody}
\emph{Self-custody}\index{self-custody} is holding crypto assets under \omterm{def:m3:blockchains-for-traders:key}{private keys} the owner controls, rather than as a balance at a venue or custodian. A
\emph{hot wallet}\index{hot wallet} is one whose keys are on internet-connected systems, able to sign at once; a \emph{cold wallet}\index{cold wallet} is one
whose keys are kept offline, so that signing requires a deliberate offline process. \emph{Multi-party computation}\index{multi-party computation} (MPC) is a
cryptographic method by which several parties, each holding a share of a key, jointly produce a signature without the key ever existing in one place.
\end{definition}

A trading firm's treasury holds most coins cold and moves only what trading needs to \omterm{def:m3:the-crypto-trading-business:custody}{hot wallets} and venues, with MPC or multi-signature approval so that no
single person or server can move funds. The failures are old ones in new dress. FTX kept almost all its crypto in \omterm{def:m3:the-crypto-trading-business:custody}{hot wallets} without multi-signature controls
(\cref{ch:m3:centralised-exchanges}). At QuadrigaCX, a Canadian exchange that collapsed in 2019, the Ontario Securities Commission's staff found that the
co-founder and chief executive had credited himself fictitious balances under aliases and traded them against clients, and had lost more on outside
platforms; over 76\,000 clients lost at least 169 million Canadian dollars in all. Staff called it ``an old-fashioned fraud wrapped in modern
technology''. Segregation of duties, reconciliation of on-chain balances against books every day, and limits on who can approve a transfer are the defences, as
in any other asset class.

\begin{omfigure}
\centering
\begin{tikzpicture}[font=\small,
  w/.style={draw,thick,rounded corners=2pt,align=center,minimum height=1cm,text width=2.5cm},
  arr/.style={-{Stealth[length=2.2mm]},thick}]
\node[w,draw=omDef,fill=omDef!8] (c) at (0,0) {cold storage\\(most assets)};
\node[w,draw=omThm,fill=omThm!8] (m) at (3.6,0) {MPC or multisig\\approvals};
\node[w,draw=omExo,fill=omExo!8] (h) at (7.2,0) {hot wallets\\(daily needs)};
\node[w,draw=omProp,fill=omProp!8] (v) at (10.8,0) {venues, OTC\\counterparties};
\draw[arr] (c) -- (m); \draw[arr] (m) -- (h); \draw[arr] (h) -- (v);
\draw[arr,dashed] (v.south) to[bend left=20] node[below,font=\footnotesize] {sweeps back of excess balances} (c.south);
\end{tikzpicture}

\omcaption{A trading firm's custody chain: assets move out of cold storage only through multi-party approval, \omterm{def:m3:the-crypto-trading-business:custody}{hot wallets} hold what the day needs, and excess
balances on venues are swept back. Schematic.}\label{fig:m3:the-crypto-trading-business:custody}
\end{omfigure}

\section{Round-the-clock operations}

Crypto markets do not close. A firm needs systems and people for every hour: monitoring of positions, margins and collateral on every venue, coverage of
weekends when banks are shut but coins move, procedures for chain upgrades and venue maintenance windows, and on-call engineers for the connectivity of
\cref{ch:m3:crypto-data-and-infrastructure}. The liquidity of a weekend is thinner and its moves larger relative to depth; the liquidations of
\cref{ch:m3:margin-liquidation-and-loss-allocation} happen when they happen. The firm's risk limits must be enforced by machines that do not sleep, with humans
alerted by exception, and its fiat liquidity must be positioned before the banks close on Friday.

\section{Regulation}

\begin{definition}[Crypto-asset service provider, Markets in Crypto-Assets Regulation]\label{def:m3:the-crypto-trading-business:mica}
A \emph{crypto-asset service provider}\index{crypto-asset service provider} (CASP) is a firm that provides crypto-asset services such as custody, operating a
trading platform, exchange of crypto assets, execution or reception and transmission of orders, for which EU law requires authorisation. The \emph{Markets in
Crypto-Assets Regulation}\index{Markets in Crypto-Assets Regulation} (MiCA) is the European Union regulation, Regulation (EU) 2023/1114, that sets rules for
issuers of crypto assets and \omterm{def:m3:blockchains-for-traders:stable}{stablecoins} and for CASPs.
\end{definition}

\begin{dated}{2026-09}{MiCA's timetable}\label{dat:m3:the-crypto-trading-business:mica}
ESMA states that MiCA entered into force in June 2023; its titles on asset-referenced \omterm{def:m3:blockchains-for-traders:key}{tokens} and e-money \omterm{def:m3:blockchains-for-traders:key}{tokens} (\omterm{def:m3:blockchains-for-traders:stable}{stablecoins}) applied from 30 June 2024 and the
rest from 30 December 2024; under its Article 143(3), entities already providing crypto-asset services may continue until 1 July 2026 or until they are granted or
refused authorisation.
\end{dated}

\begin{definition}[Travel rule]\label{def:m3:the-crypto-trading-business:travel}
The \emph{travel rule}\index{travel rule} is the requirement, in the FATF's Recommendation 16 on wire transfers, that information about the originator and the
beneficiary accompany a transfer between regulated firms; the FATF extended its standards to virtual assets and their service providers in 2018.
\end{definition}

The FATF's interpretive note of June 2019 applies Recommendation 16 to virtual-asset transfers, and lets countries exempt transfers below a threshold
no higher than USD/EUR~1\,000. The European Union chose none: its Regulation 2023/1113, applying from 30 December 2024, makes the originator's
name, ledger address and identifying details travel with every transfer of crypto-assets, whatever its amount.

\begin{dated}{2026-09}{The US stablecoin law}\label{dat:m3:the-crypto-trading-business:genius}
The GENIUS Act, Public Law 119-27 of 18 July 2025, provides for the regulation of payment \omterm{def:m3:blockchains-for-traders:stable}{stablecoins}; it requires a permitted issuer to maintain identifiable
reserves on an at least one-to-one basis, in coins and currency, demand deposits, Treasury bills with a maturity of 93 days or less, repurchase agreements and
similar assets.
\end{dated}

For a trading firm, regulation changes the map of \cref{ch:m3:centralised-exchanges}: which venues it may use and from where, whether its own activities (custody
for clients, market making for issuers, OTC trading) require authorisation, what it must report and to whom, and what information must travel with its
transfers.

\section{Tutorial: valuing a listing deal}\label{tut:m3:the-crypto-trading-business:deal}

\begin{tutorial}
\textbf{Goal.} Value a \omterm{def:m3:blockchains-for-traders:key}{token} loan with three call strikes by simulation, derive the fee it implies, and compute the hedge the market maker must hold at launch
and through the year. \textbf{End state:} \cref{fig:m3:the-crypto-trading-business:fee,fig:m3:the-crypto-trading-business:hedge} and the numbers of the
weekend problem.
\begin{tutsteps}
\item \textbf{The package.} Monte Carlo with antithetic pairs and the terminal price as a control variate, checked against Black.
\omcode{code/firm/tokenloan/firm_tokenloan.py}{39}{55}{Monte Carlo value of the call package with a control variate.}\label{lst:m3:the-crypto-trading-business:mc}
\item \textbf{Fee and hedge.} \texttt{implied\_fee}, \texttt{day\_one\_hedge} and \texttt{delta\_schedule}.
\omcode{code/firm/tokenloan/firm_tokenloan.py}{62}{74}{The implied fee, the day-one hedge and the hedge through the year.}\label{lst:m3:the-crypto-trading-business:hedge}
\item \textbf{Run} \texttt{summary()} and \texttt{fig\_business.py}.
\end{tutsteps}
\textbf{What to change next.} Add a \omterm{def:m3:the-crypto-trading-business:unlock}{token unlock} of 10\% of supply at month six as a jump in supply that moves the price; charge the hedge the spread it pays
on a thin market; value the deal from the issuer's side, including the \omterm{def:m3:blockchains-for-traders:key}{tokens} it gets back.
\end{tutorial}

\section{Build: token market-making agreements}\label{bld:m3:the-crypto-trading-business:tokenloan}

\begin{build}
\bfield{Purpose} The miniature firm signs, values and hedges \omterm{def:m3:the-crypto-trading-business:tmma}{token market-making agreements}, monitors its quoting obligations, and books the value it earns.
\bfield{Interface} \texttt{Deal(tokens\_loaned, price, years, vol, tranches)}; \texttt{call\_bs}; \texttt{package\_value\_mc}, \texttt{package\_value};
\texttt{implied\_fee}; \texttt{day\_one\_hedge}; \texttt{delta\_schedule(deal, prices, t\_left)}; \texttt{obligations\_met(spread\_bp, depth\_usd, uptime)};
\texttt{post\_fee(ledger, deal, ts)} posting to \texttt{firm.ledger}.
\bfield{Rules} \omterm{def:m3:blockchains-for-traders:key}{Tokens} received and owed netted; volatility from the \omterm{def:m3:blockchains-for-traders:key}{token}'s own history and comparable listings, never assumed low; obligations checked per venue
and per day; the calls booked at signing as fee revenue in the ledger.
\bfield{Acceptance tests} \texttt{code/firm/tokenloan/tests/}: Monte Carlo within 1\% of the closed form; call limits; fee between 0 and 1; a hedge that rises with the
price; obligations; a ledger posting.
\bfield{Stretch} Unlock schedules; hedging costs on thin books; accounting for the loan liability; the deal's value to the issuer.
\end{build}

\begin{omsources}
\item SEC, press release 2024-166, 9 October 2024.
\item ESMA, Markets in Crypto-Assets Regulation (MiCA) page; Regulation (EU) 2023/1114. Accessed September 2026.
\item GENIUS Act, Public Law 119-27, 18 July 2025 (GovInfo).
\item Ontario Securities Commission, \emph{QuadrigaCX: A Review by Staff of the Ontario Securities Commission}, 2020.
\item FATF, \emph{International Standards on Combating Money Laundering and the Financing of Terrorism and Proliferation} (the FATF Recommendations),
updated June 2025: Recommendations 15 and 16 and the Interpretive Note to Recommendation 15 (Internet Archive copy).
\item Regulation (EU) 2023/1113 of 31 May 2023 on information accompanying transfers of funds and certain crypto-assets (Internet Archive copy of EUR-Lex).
\item US Attorney's Office, District of Massachusetts, press release of 9 October 2024 (Internet Archive copy).
\end{omsources}

\section{Exercises}

\begin{exercise}[$\star$]\label{exo:m3:the-crypto-trading-business:1}
Why does a \omterm{def:m3:blockchains-for-traders:key}{token} loan by itself leave the market maker with no net exposure to the \omterm{def:m3:blockchains-for-traders:key}{token}?
\end{exercise}

\begin{exercise}[$\star$]\label{exo:m3:the-crypto-trading-business:2}
Distinguish a \omterm{def:m3:the-crypto-trading-business:custody}{hot wallet} from a \omterm{def:m3:the-crypto-trading-business:custody}{cold wallet}, and say what MPC adds.
\end{exercise}

\begin{exercise}[$\star$]\label{exo:m3:the-crypto-trading-business:3}
From when did MiCA apply to \omterm{def:m3:blockchains-for-traders:stable}{stablecoin} issuers, and from when to \omterm{def:m3:the-crypto-trading-business:mica}{crypto-asset service providers}?
\end{exercise}

\begin{exercise}[$\star\star$]\label{exo:m3:the-crypto-trading-business:4}
Value one call on a \omterm{def:m3:blockchains-for-traders:key}{token} at USD~0.50, strike 0.75, one year, volatility 120\%, without drift or discounting.
\end{exercise}

\begin{exercise}[$\star\star$]\label{exo:m3:the-crypto-trading-business:5}
How does the implied fee of the listing deal change between volatilities of 80\% and 160\%? Why so much?
\end{exercise}

\begin{exercise}[$\star\star$]\label{exo:m3:the-crypto-trading-business:6}
An OTC client asks for a price on 500 bitcoin. List what goes into the desk's quote.
\end{exercise}

\begin{exercise}[$\star\star\star$]\label{exo:m3:the-crypto-trading-business:7}
\emph{Coding.} With \texttt{delta\_schedule}, how many \omterm{def:m3:blockchains-for-traders:key}{tokens} must the market maker buy back if the price falls from 0.50 to 0.25 with half a year left?
\end{exercise}

\begin{exercise}[$\star\star\star$]\label{exo:m3:the-crypto-trading-business:8}
\emph{Find the flaw.} ``Our agreement pays us only if volume on the \omterm{def:m3:blockchains-for-traders:key}{token} reaches USD~5 million a day; we will make sure it does.''
\end{exercise}

\section{Problem: The Listing Deal}

\begin{problem}[{Weekend problem --- what the calls are worth and what to sell}]\label{pb:m3:the-crypto-trading-business:1}
A project with 100 million \omterm{def:m3:blockchains-for-traders:key}{tokens} lists at USD~0.50. It lends a market maker 2\% of the supply for a year and grants calls on a third of that amount each at
USD~0.75, 1.00 and 1.50. The \omterm{def:m3:blockchains-for-traders:key}{token}'s volatility is estimated at 120\% a year; ignore interest and drift.

\medskip\noindent\textbf{Part I --- The value.}
\begin{enumerate}
\item What is the loan worth in dollars?
\item What are the three call tranches worth, and the package?
\item What fee does the package imply, as a share of the loan?
\item How close does the Monte Carlo value come to the closed form?
\item At 80\% volatility, what is the fee?
\end{enumerate}

\medskip\noindent\textbf{Part II --- The hedge.}
\begin{enumerate}[resume]
\item How many \omterm{def:m3:blockchains-for-traders:key}{tokens} does the market maker sell on day one?
\item What share of the loan is that?
\item What happens to the hedge if the \omterm{def:m3:blockchains-for-traders:key}{token} doubles in the first month?
\item And if it halves?
\item Why does the day-one sale matter to the new market?
\end{enumerate}

\medskip\noindent\textbf{Part III --- The obligations.}
\begin{enumerate}[resume]
\item What obligations should the agreement set, and which should it avoid?
\item How would the market maker meet a maximum spread of 1\% with depth of USD~50\,000 on each side?
\item What does a \omterm{def:m3:the-crypto-trading-business:unlock}{token unlock} at month six do to the hedge and the quotes?
\item How should the firm book the deal?
\item What would make this deal unlawful?
\end{enumerate}

\medskip\noindent\textbf{Part IV --- Judgement.}
\begin{enumerate}[resume]
\item Who has the better side of the deal?
\item Why do issuers pay in calls rather than cash?
\item How would you check that a market maker is honest?
\item State the \emph{named result}: the value of the call package as a share of the loaned \omterm{def:m3:blockchains-for-traders:key}{tokens}' value, and the \omterm{def:m3:blockchains-for-traders:key}{tokens} sold on day one.
\item In one sentence: what does a listing market maker sell to the issuer?
\end{enumerate}
\end{problem}

\section{Interview questions}

\begin{interviewq}[$\star$ \iqroles{trader}]\label{iq:m3:the-crypto-trading-business:1}
How is a market maker paid in a loan-plus-call deal, and how does it hedge?
\end{interviewq}

\begin{interviewq}[$\star$ \iqroles{risk}]\label{iq:m3:the-crypto-trading-business:2}
Describe a sound custody set-up for a crypto trading firm.
\end{interviewq}

\begin{interviewq}[$\star\star$ \iqroles{researcher}]\label{iq:m3:the-crypto-trading-business:3}
How would you estimate the volatility of a \omterm{def:m3:blockchains-for-traders:key}{token} that has not traded yet?
\end{interviewq}

\begin{interviewq}[$\star\star$ \iqroles{trader}]\label{iq:m3:the-crypto-trading-business:4}
How do you price an OTC block of 500 bitcoin at 3 a.m. on a Sunday?
\end{interviewq}

\begin{interviewq}[$\star\star$ \iqroles{risk}]\label{iq:m3:the-crypto-trading-business:5}
What changes for a trading firm when MiCA and the \omterm{def:m3:the-crypto-trading-business:travel}{travel rule} apply to it?
\end{interviewq}

\begin{interviewq}[$\star\star\star$ \iqroles{developer, risk}]\label{iq:m3:the-crypto-trading-business:6}
Design the controls for moving USD~50 million of coins from cold storage to three venues within an hour.
\end{interviewq}
